\subsection{Prediction: $G_{\mathrm{eff}}$ in terms of $c_\chi$ and $\hbar_\chi$}
\label{subsec:stepB-G}

Comparing~\eqref{eq:BI-IR} with the standard Einstein--Hilbert normalization
$\frac{c^4}{16\pi G} \int R \sqrt{-g}\, d^4x$ and using $c=c_\chi$, the identification
requires:
\[
  \frac{c_\chi^4}{16\pi G_{\mathrm{eff}}} = \frac{c_\chi^2}{2\hbar_\chi}.
\]
Formally, this gives:
\begin{equation}
  G_{\mathrm{eff}} = \frac{c_\chi^2 \hbar_\chi}{8\pi}.
  \label{eq:G-prediction}
\end{equation}

A dimensional check in SI units yields $[c_\chi^2 \hbar_\chi] = \mathrm{kg}\,\mathrm{m}^4\,\mathrm{s}^{-3}$, while
$[G]=\mathrm{m}^3\,\mathrm{kg}^{-1}\,\mathrm{s}^{-2}$.
These do not match, which signals that the prefactor $c_\chi^4/\hbar_\chi^2$ in
Eq.~\eqref{eq:BI-gravity} must carry the residual dimensions.
The dimensional consistency of Eq.~\eqref{eq:G-prediction} depends on the precise
normalization of the Born--Infeld gravitational action, which requires the full
tensorial derivation noted in Section~\ref{subsec:stepB-BI} as an open step.
The structural prediction $G_{\mathrm{eff}} = G(c_\chi, \hbar_\chi)$ is robust.
The precise numerical prefactor and dimensional matching are left as a quantitative
step contingent on the tensorial completion.
